\section{Modelos de difusión discreta: adición y eliminación de ruido en secuencias de proteínas}

\subsection{Estrategias de adición de ruido para datos discretos}

Para datos de secuencias discretas (como el lenguaje natural o las secuencias de proteínas), Amini propuso dos estrategias para implementar la «adición de ruido»:

\textbf{Estrategia uno: Masking (enmascaramiento)}

Partiendo de una secuencia limpia, en cada paso se elige aleatoriamente una o varias posiciones y se sustituye el token correspondiente por la marca especial \texttt{[MASK]}:

\begin{table}[H]
\centering
\caption{Ilustración de la estrategia Masking (con el ejemplo en lenguaje natural «This class is fun»)}
\begin{tabular}{@{}cl@{}}
\toprule
\textbf{Paso} & \textbf{Estado de la secuencia} \\
\midrule
$t=0$ & This class is fun \\
$t=1$ & This {[}MASK{]} is fun \\
$t=2$ & {[}MASK{]} {[}MASK{]} is fun \\
$t=3$ & {[}MASK{]} {[}MASK{]} {[}MASK{]} fun \\
$t=4$ & {[}MASK{]} {[}MASK{]} {[}MASK{]} {[}MASK{]} \\
\bottomrule
\end{tabular}
\end{table}

\textbf{Estrategia dos: Mutation (mutación/sustitución)}

No se usa una marca especial; en cambio, en cada paso, con cierta probabilidad, se sustituyen aleatoriamente los tokens de algunas posiciones por otros tokens del vocabulario:

\begin{table}[H]
\centering
\caption{Ilustración de la estrategia Mutation}
\begin{tabular}{@{}cl@{}}
\toprule
\textbf{Paso} & \textbf{Estado de la secuencia} \\
\midrule
$t=0$ & This class is fun \\
$t=1$ & This \textit{apple} is fun \\
$t=2$ & \textit{Rock} \textit{apple} is fun \\
$t=3$ & \textit{Rock} \textit{apple} \textit{my} fun \\
$t=T$ & (secuencia de tokens completamente aleatoria) \\
\bottomrule
\end{tabular}
\end{table}

\begin{warningbox}{Mutation es más difícil que Masking}
La estrategia Mutation exige más del modelo de eliminación de ruido, porque este no puede identificar mediante una marca especial qué posiciones fueron modificadas: tiene que determinar por sí mismo qué tokens son «incorrectos» y predecir el valor correcto. Esto, en cierto sentido, se acerca más al proceso real de mutación evolutiva, pero también aumenta considerablemente la dificultad del entrenamiento. En EvoDiff se adoptaron y probaron ambas estrategias.
\end{warningbox}

\subsection{Formalización de la difusión discreta}

Para una secuencia de proteína $\mathbf{s} = (s_1, s_2, \ldots, s_L)$, donde $s_i \in \{A_1, A_2, \ldots, A_{20}\}$ representa uno de los 20 aminoácidos:

\textbf{Proceso hacia adelante}: en cada paso $t$, se elige aleatoriamente un subconjunto de posiciones $\mathcal{M}_t$ y se sustituyen los aminoácidos de esas posiciones por \texttt{[MASK]}:

$$s_i^{(t)} = \begin{cases} \texttt{[MASK]} & \text{if } i \in \mathcal{M}_1 \cup \mathcal{M}_2 \cup \cdots \cup \mathcal{M}_t \\ s_i^{(0)} & \text{otherwise} \end{cases}$$

\textbf{Proceso inverso}: se entrena una red neuronal $f_\theta$ que, en cada paso, restaura algunas de las posiciones enmascaradas a su aminoácido correcto:

$$f_\theta(\mathbf{s}^{(t)}, t) \rightarrow \mathbf{s}^{(t-1)}$$

Al generar, se parte de la secuencia totalmente enmascarada $\mathbf{s}^{(T)} = (\texttt{[MASK]}, \texttt{[MASK]}, \ldots, \texttt{[MASK]})$ y se elimina el ruido paso a paso hasta obtener una secuencia de proteína completa.

\subsection{Relación entre la difusión discreta y los modelos de lenguaje}

Amini señaló una conexión teórica profunda: \textbf{la difusión discreta con Masking es, en realidad, una generalización de varios esquemas de modelado de lenguaje}:

\begin{table}[H]
\centering
\caption{Comparación entre la difusión discreta y los esquemas clásicos de modelado de lenguaje}
\begin{tabular}{@{}lll@{}}
\toprule
\textbf{Método} & \textbf{Orden de decodificación} & \textbf{Predicción por paso} \\
\midrule
Next Token Prediction (LLM) & Fijo, de izquierda a derecha & Un solo token siguiente \\
Masked Language Model (BERT) & Toda la secuencia visible & Se rellenan todos los mask en un paso \\
Discrete Diffusion & Orden aleatorio, paso a paso & Se rellena parte de los mask en cada paso \\
\bottomrule
\end{tabular}
\end{table}

\begin{importantbox}{Capacidad de generalización de la difusión discreta}
El modelo de difusión discreta aprende la tarea de eliminación de ruido sobre \textbf{todos los órdenes de decodificación} y \textbf{todos los subconjuntos de masking} posibles. Esto hace que abarque, al mismo tiempo:
\begin{itemize}
    \item la decodificación «de izquierda a derecha» de los modelos autorregresivos (un orden de masking particular)
    \item la «eliminación de ruido en un paso» de Masked Language Model (una configuración particular del número de pasos)
\end{itemize}
Por lo tanto, el modelo de difusión discreta es un marco más general, con mayor flexibilidad y capacidad expresiva.
\end{importantbox}

\subsection{Resumen de la sección}

Trasladar el Diffusion Model del dominio continuo al dominio discreto requiere redefinir el «ruido»: mediante Masking o Mutation se corrompen progresivamente los tokens de la secuencia. El marco de difusión discreta unifica matemáticamente los modelos autorregresivos y Masked Language Model, y ofrece una herramienta poderosa y flexible para el modelado de secuencias de proteínas.


